\chapter{Бимодульный селектор кратности и расширение полного поля}
\label{ch:v8-bimodule-multiplicity-separator-gate}

\section{Пятнадцатимерный модуль недостаточен}

Полный калибровочно-замкнутый модуль переноса предыдущих гейтов имел размерность
15 над $\mathbb C$. Но физическая конечная алгебра различает концевые блоки
\begin{equation}
 E_s=Q_L\oplus L_L\oplus X_L\oplus Y_L,
 \qquad
 E_t=u_R\oplus d_R\oplus e_R\oplus X_R\oplus Y_R.
 \label{eq:v8-endpoint-bimodule-blocks}
\end{equation}
Для центральных проекторов $P_s^a,P_t^b$ действие на стрелке равно
\begin{equation}
 A\longmapsto P_t^bAP_s^a.
 \label{eq:v8-endpoint-block-action}
\end{equation}
Старое пространство не замкнуто относительно этих операций: максимальный
остаток проекции назад в 15-мерный модуль равен $0.429699\ldots$.

Минимальное замыкание под всеми концевыми проекторами имеет размерность
\begin{equation}
 \dim_{\mathbb C}\mathcal T_{\rm bimod}=20.
 \label{eq:v8-bimodule-transfer-dimension}
\end{equation}
Оно автоматически калибровочно-замкнуто: последовательность замыкания равна
$20\to20$. Новых концевых фермионов и калибровочных генераторов не добавляется, но
появляются пять дополнительных комплексных компонент скалярного поля переноса.
Полное внутреннее пространство полей теперь имеет
\begin{equation}
 2\cdot20+12=52
 \label{eq:v8-bimodule-full-field-real-dimension}
\end{equation}
вещественные компоненты.

\section{Структурное разделение инцидентностной и тяжёлой частей}

На бимодульно замкнутом носителе опора инцидентности состоит из трёх
исходных пар $H_{15}$ и четырёх изотипических пар:
\begin{align}
 E_I=\{&u_RQ_L,d_RQ_L,e_RL_L,
 e_RX_L,X_RX_L,Y_RL_L,Y_RY_L\}.
 \label{eq:v8-bimodule-incidence-support}
\end{align}
Оставшаяся тяжёлая опора равна
\begin{equation}
 E_H=\{d_RX_L,e_RY_L,X_RL_L,Y_RQ_L\}.
 \label{eq:v8-bimodule-heavy-support}
\end{equation}
Эти множества определены концевыми метками и равенством калибровочных типов, а не
амплитудами выбранного вакуума.

Соответствующие подмодули ортогональны и имеют одинаковую размерность:
\begin{equation}
 \mathcal T_{\rm bimod}=\mathcal I_{10}\oplus\mathcal H_{10}.
 \label{eq:v8-bimodule-ten-plus-ten}
\end{equation}
Сумма блоковых проекторов на $E_I$ воспроизводит проектор инцидентности с
остатком меньше $10^{-15}$. Тем самым кратность $4+4$, невидимая для калибровочной и
данным вещественной структуры на сжатом модуле, полностью разделяется бимодульными метками.

\section{Единый Ходж-уровень}

Общий концевой след даёт одинаковый вес каждой компоненте переноса. Так как
оба подмодуля имеют размерность 10 над $\mathbb C$, их полные следовые веса
равны при любой центральной концевой плотности. Один Ходж-уровень задаёт
потенциал
\begin{equation}
 S_E=\sum_{a\in E_I}(|z_a|^2-1)^2
 +\sum_{b\in E_H}\bigl[(|w_b|^2+1)^2-1\bigr].
 \label{eq:v8-bimodule-common-level-potential}
\end{equation}
Его гессиан имеет массы $-4,+4$ в нуле и $8,+4$ в вакууме. Поэтому прежняя
свобода $m_I/m_H$ устраняется внутри рёберно-ходжева блока.

После добавления относительного Грам-гессиана с весом $\beta$ и гладкой
изотипической кривизны точное окно запуска равно
\begin{equation}
 0\leq\beta<\frac23.
 \label{eq:v8-bimodule-beta-window}
\end{equation}
При $\beta=1/2$ получено
\begin{equation}
 (n_-,n_0,n_+)_{0}=(20,0,20),
 \qquad
 (n_-,n_0,n_+)_{*}=(0,0,40),
 \qquad
 \lambda_{\min}^{*}=5.773318\ldots.
 \label{eq:v8-bimodule-full-transition}
\end{equation}
На границе $\beta=2/3$ появляются двенадцать нулевых мод, а при $\beta=1$
исходная сигнатура становится $(38,0,2)$.

\begin{theorem}[Бимодульное снятие кратности]
Минимальное замыкание поля переноса относительно полных концевых проекторов
имеет размерность 20 над $\mathbb C$ и канонически раскладывается как
$\mathcal I_{10}\oplus\mathcal H_{10}$. Один общий Ходж-уровень фиксирует
равные рёберные массы и даёт качественный переход
$(20,0,20)\to(0,0,40)$ во всём окне $0\leq\beta<2/3$.
\end{theorem}

\begin{nogocriterion}[Относительный Грам-вес остаётся открытым]
Нельзя считать полный родительское действие полученным: бимодульная типизация
устраняет отношение $m_I/m_H$, но не выводит вес $\beta$ между рёберно-ходжева и
грамова метрика концевых секторов кривизнами. Общий равный вес $\beta=1$ не проходит исходный
тест. Кроме того, расширение $15\to20$ должно быть встроено в одну
суперсвязностную кривизну, а не только в её гессиан.
\end{nogocriterion}

\begin{protocol}
\begin{enumerate}
 \item старый модуль переноса: \textbf{$15$ комплексных измерений};
 \item его бимодульная замкнутость: \textbf{нет};
 \item минимальное замыкание: \textbf{$20$ комплексных измерений, $40$ вещественных измерений};
 \item полное внутреннее пространство полей: \textbf{$52$ вещественных измерения};
 \item инцидентностная и тяжёлая части: \textbf{$10+10$ комплексных измерений};
 \item новых концевых фермионов: \textbf{$0$};
 \item новых калибровочных генераторов: \textbf{$0$};
 \item бимодульный проектор инцидентности: \textbf{получен};
 \item отношение рёберных масс: \textbf{$1$ из общего уровня};
 \item окно $\beta$: \textbf{$0\leq\beta<2/3$};
 \item переход при $\beta=1/2$: \textbf{$(20,0,20)\to(0,0,40)$};
 \item происхождение $\beta$: \textbf{открыто};
 \item полный родительское действие: \textbf{не получен};
 \item следующий гейт: \textbf{общая кривизна и относительный вес}.
\end{enumerate}
\end{protocol}

Машинный сертификат сохранён в
\path{s2t_v8_bimodule_multiplicity_separator_gate_results.json}.